% ============================================================
% CHAPITRE 1 — ÉTUDE DU PROJET
% ============================================================

\chapter{Étude du Projet}

\section*{Plan}

\begin{tabular}{r l r}
\textbf{1} & \textbf{Organisme d'accueil} & \textbf{\pageref{sec:organisme}} \\
\textbf{2} & \textbf{Présentation du projet} & \textbf{\pageref{sec:projet_overview}} \\
\textbf{3} & \textbf{Méthodologie de développement} & \textbf{\pageref{sec:methodologie}} \\
\textbf{4} & \textbf{Méthodologies d'ingénierie complémentaires} & \textbf{\pageref{sec:methodologies_complementaires}} \\
\textbf{5} & \textbf{Langage de modélisation~: UML} & \textbf{\pageref{sec:uml}} \\
\end{tabular}

\section*{Introduction}

Ce chapitre introductif pose les fondations de notre projet de fin d'études. Nous débutons par une présentation détaillée de l'organisme d'accueil, \textbf{TIM Tunisie}, en décrivant sa structure organisationnelle et ses produits phares. Par la suite, nous exposons la problématique qui motive ce travail, la solution proposée pour y répondre ainsi que les objectifs visés. Enfin, nous décrivons la méthodologie de développement retenue pour conduire ce projet dans un cadre structuré et itératif.

% ────────────────────────────────────────────────────────────
\section{Organisme d'accueil}
\label{sec:organisme}
% ────────────────────────────────────────────────────────────

\subsection{TIM Tunisie --- Présentation générale}

\textbf{TIM Tunisie} (\textit{Technologies, Informatique et Management}) est une entreprise tunisienne spécialisée dans la conception et le développement de solutions logicielles à destination du secteur industriel. Fondée avec pour vocation d'accompagner la transformation numérique des entreprises industrielles en Tunisie et dans la région MENA, TIM Tunisie a acquis une expertise reconnue dans deux domaines stratégiques~:

\begin{itemize}[noitemsep]
    \item la \textbf{gestion de la qualité} industrielle conformément aux normes internationales~;
    \item la \textbf{gestion de la maintenance} assistée par ordinateur (GMAO).
\end{itemize}

L'entreprise se positionne aujourd'hui comme un acteur clé de l'innovation technologique industrielle, en intégrant progressivement des composantes d'intelligence artificielle --- et notamment des modèles de langage de grande taille --- au sein de ses produits existants.

\begin{figure}[H]
    \centering
    \fbox{\parbox{6cm}{\centering\vspace{1cm}\textit{Logo de TIM Tunisie}\vspace{1cm}}}
    \caption{Logo de TIM Tunisie}
    \label{fig:logo_tim}
\end{figure}

\begin{table}[H]
\centering
\caption{Fiche signalétique de TIM Tunisie}
\label{tab:fiche_tim}
\begin{tabularx}{\textwidth}{|l|X|}
\hline
\textbf{Raison sociale} & TIM Tunisie \\
\hline
\textbf{Secteur d'activité} & Édition de logiciels industriels \\
\hline
\textbf{Produits principaux} & QALITAS, GMAO PRO \\
\hline
\textbf{Domaines} & Qualité industrielle, Maintenance, Intelligence Artificielle \\
\hline
\textbf{Site web} & \url{https://tim-tunisie.com/} \\
\hline
\textbf{Localisation} & Tunis, Tunisie \\
\hline
\end{tabularx}
\end{table}


\subsection{Les produits de TIM Tunisie}

L'activité de TIM Tunisie s'articule autour de deux produits logiciels complémentaires, chacun répondant à des besoins industriels spécifiques et intégrant des fonctionnalités d'intelligence artificielle.

\subsubsection{QALITAS --- Plateforme de gestion de la qualité}

\textbf{QALITAS} est un système d'information intégré dédié à la gestion de la qualité industrielle. Cette plateforme permet aux entreprises de structurer et d'automatiser l'ensemble de leurs processus qualité~:

\begin{itemize}[noitemsep]
    \item gestion des processus de contrôle qualité et d'audit interne~;
    \item suivi des non-conformités et des actions correctives et préventives~;
    \item production de rapports de conformité aux normes internationales (ISO~9001, ISO~14001, ISO~45001)~;
    \item automatisation des workflows de validation et d'approbation documentaire.
\end{itemize}

Avec l'intégration progressive de l'intelligence artificielle, QALITAS exploite désormais des \textbf{modèles de langage (LLM)} pour l'analyse automatique de documents qualité, la génération de rapports de conformité et l'assistance conversationnelle pour les auditeurs internes.

\subsubsection{GMAO PRO --- Gestion de la maintenance assistée par ordinateur}

\textbf{GMAO PRO} constitue la solution de TIM Tunisie pour la gestion de la maintenance industrielle. Ses fonctionnalités principales couvrent~:

\begin{itemize}[noitemsep]
    \item la planification et le suivi des opérations de maintenance préventive et corrective~;
    \item la gestion du parc d'équipements et des pièces de rechange~;
    \item le suivi des indicateurs clés de performance de maintenance (MTBF, MTTR, taux de disponibilité)~;
    \item la génération automatisée de bons de travail et de rapports d'intervention.
\end{itemize}

GMAO PRO intègre également des fonctionnalités d'IA pour la \textbf{maintenance prédictive}, s'appuyant sur les LLM pour analyser les historiques de pannes et formuler des recommandations d'intervention anticipée.


% ────────────────────────────────────────────────────────────
\section{Présentation du projet}
\label{sec:projet_overview}
% ────────────────────────────────────────────────────────────

Ce projet s'inscrit dans le cadre d'un stage de fin d'études en licence en Informatique de Gestion à l'Institut Supérieur de Gestion de Tunis. Le stage a été effectué au sein de TIM Tunisie sous la supervision académique de \textbf{Mme Salsabil Hadji}.

\subsection{Étude de l'existant}

Avant le développement de notre plateforme, l'utilisation des LLM au sein des produits de TIM Tunisie se caractérisait par plusieurs lacunes structurelles~:

\begin{itemize}[noitemsep]
    \item \textbf{Appels directs}~: chaque module applicatif (QALITAS, GMAO PRO) effectuait des appels directs aux API des fournisseurs LLM (OpenAI, Anthropic, Google) sans intermédiaire centralisé~;
    \item \textbf{Absence de proxy}~: aucun point de passage unique ne permettait d'intercepter, de monitorer ou de gouverner les appels~;
    \item \textbf{Clés API partagées}~: les clés d'accès étaient souvent partagées entre plusieurs services, sans isolation par client~;
    \item \textbf{Pas de comptage de tokens}~: la consommation réelle n'était pas tracée de manière systématique~;
    \item \textbf{Facturation manuelle}~: les coûts LLM étaient estimés manuellement à partir des factures des fournisseurs.
\end{itemize}

\subsection{Critique de l'existant}

L'analyse approfondie de la situation existante a révélé des \textbf{lacunes critiques} compromettant la viabilité économique et opérationnelle de l'utilisation des LLM~:

\begin{table}[H]
\centering
\caption{Critique de l'existant --- Problèmes identifiés}
\label{tab:critique_existant}
\begin{tabularx}{\textwidth}{|c|p{4.5cm}|X|}
\hline
\textbf{\#} & \textbf{Problème} & \textbf{Impact} \\
\hline
1 & Pas de visibilité sur la consommation & Impossible de savoir quel agent, modèle ou client consomme le plus \\
\hline
2 & Pas de contrôle des coûts & Risque de dépassements budgétaires non détectés \\
\hline
3 & Pas d'isolation multi-tenant & Un client peut monopoliser les ressources sans limitation \\
\hline
4 & Pas de gouvernance & Aucune règle de quota, de budget ou de restriction de modèle \\
\hline
5 & Pas de détection d'anomalies & Les pics de consommation anormaux passent inaperçus \\
\hline
6 & Pas de prévision budgétaire & Impossible d'anticiper les coûts des mois suivants \\
\hline
7 & Pas de facturation automatisée & Les factures produites manuellement sont sources d'erreurs \\
\hline
8 & Pas d'observabilité & Aucune métrique, aucun tableau de bord pour les décideurs \\
\hline
\end{tabularx}
\end{table}


\subsection{Problématique}

La problématique centrale de ce projet peut être formulée ainsi~:

\begin{quote}
\textit{\textbf{Comment concevoir et implémenter une plateforme centralisée capable de gouverner, monitorer, analyser et optimiser l'utilisation des modèles de langage (LLM) dans un contexte multi-tenant industriel, tout en assurant la traçabilité complète, la maîtrise des coûts et la détection proactive des anomalies de consommation~?}}
\end{quote}

Cette problématique se décline en sous-questions de recherche~:

\begin{enumerate}[noitemsep]
    \item Comment intercepter et journaliser tous les appels LLM de manière transparente pour les applications clientes~?
    \item Comment appliquer des règles de gouvernance (quotas, budgets, restrictions de modèles) en temps réel~?
    \item Comment détecter automatiquement les anomalies de consommation à l'aide d'algorithmes d'apprentissage automatique~?
    \item Comment fournir des explications intelligibles des anomalies détectées via l'IA~?
    \item Comment prédire les coûts futurs et alerter sur les risques de dépassement budgétaire~?
    \item Comment automatiser la facturation par client avec des contrats flexibles~?
    \item Comment offrir un tableau de bord exécutif interactif et en temps réel~?
\end{enumerate}


\subsection{Solution proposée}

Pour répondre à la problématique identifiée, nous proposons le développement de \textbf{LLM Dashboard}, une plateforme SaaS de gouvernance et de monitoring des modèles de langage. Cette solution s'articule autour de \textbf{onze modules fonctionnels} interconnectés, formant un écosystème cohérent pour la gestion complète du cycle de vie des appels LLM~:

\begin{table}[H]
\centering
\caption{Modules fonctionnels de la plateforme LLM Dashboard}
\label{tab:modules}
\begin{tabularx}{\textwidth}{|c|l|X|}
\hline
\textbf{N°} & \textbf{Module} & \textbf{Description} \\
\hline
M1 & Gateway & Proxy LLM intelligent avec routage, gouvernance temps réel et journalisation complète \\
\hline
M2 & Analytics & Agrégation des coûts quotidiens et mensuels, pricing versionné par modèle \\
\hline
M3 & Detection & Détection d'anomalies par ML (STL, Isolation Forest, CUSUM) \\
\hline
M4 & Explainer & Explication automatique des anomalies détectées via LLM multi-fournisseur \\
\hline
M5 & Optimizer & Recommandations d'optimisation des prompts et des modèles avec simulation ROI \\
\hline
M6 & Forecasting & Prévisions budgétaires à trois scénarios et analyse de risques \\
\hline
M7 & Auth & Authentification JWT, RBAC à 4 rôles, gestion des clés API \\
\hline
M8 & Dashboard & Tableau de bord exécutif avec WebSocket temps réel et assistant IA \\
\hline
M9 & Tracing & Suivi des sessions conversationnelles, A/B testing, alertes multi-canal \\
\hline
M10 & Billing & Facturation contractuelle, génération PDF, webhooks de notification \\
\hline
M11 & Prompts & CMS de templates de prompts versionné avec variables dynamiques \\
\hline
\end{tabularx}
\end{table}


\subsection{Objectifs du projet}

Les objectifs de ce projet sont définis selon une hiérarchie de priorités~:

\begin{enumerate}[noitemsep]
    \item \textbf{Centraliser} tous les appels LLM à travers un proxy intelligent unique~;
    \item \textbf{Journaliser} chaque interaction avec les modèles (tokens, coûts, latence, statut) dans une base de données optimisée pour les séries temporelles~;
    \item \textbf{Gouverner} l'utilisation des LLM via un moteur de règles configurables~;
    \item \textbf{Détecter} automatiquement les anomalies de consommation grâce à trois algorithmes ML complémentaires~;
    \item \textbf{Expliquer} les anomalies détectées de manière intelligible via un module d'IA dédié~;
    \item \textbf{Optimiser} les coûts en recommandant des modèles alternatifs et des restructurations de prompts~;
    \item \textbf{Prévoir} les budgets futurs avec trois scénarios (optimiste, moyen, pessimiste)~;
    \item \textbf{Facturer} automatiquement chaque tenant selon son type de contrat~;
    \item \textbf{Visualiser} toutes les métriques dans un tableau de bord interactif~;
    \item \textbf{Sécuriser} la plateforme avec un système d'authentification robuste et une isolation stricte par tenant.
\end{enumerate}


% ────────────────────────────────────────────────────────────
\section{Méthodologie de développement}
\label{sec:methodologie}
% ────────────────────────────────────────────────────────────

Le choix de la méthodologie de développement constitue une décision structurante pour la réussite d'un projet logiciel. Cette section présente la méthodologie \textbf{Scrum} que nous avons adoptée pour conduire ce projet, en justifiant ce choix par rapport aux alternatives existantes.

\subsection{Méthodologies agiles --- Contexte}

Les méthodologies agiles sont nées en réponse aux limites des approches traditionnelles de développement logiciel, dites \og en cascade \fg{} (\textit{waterfall}). Le \textit{Manifesto for Agile Software Development}, publié en 2001, a posé les fondements d'une nouvelle philosophie de développement centrée sur quatre valeurs fondamentales~: les individus et les interactions, les logiciels fonctionnels, la collaboration avec le client et la réponse au changement.

Parmi les frameworks agiles les plus répandus, on distingue Scrum, Kanban, Extreme Programming (XP) et SAFe. Pour notre projet, nous avons retenu \textbf{Scrum} en raison de sa structure claire, de son orientation vers les livrables incrémentaux et de sa capacité à encadrer un projet de taille moyenne avec une équipe restreinte.

\subsection{Le framework Scrum}

Scrum est un cadre de travail agile conçu pour le développement itératif et incrémental de produits complexes. Défini formellement par Schwaber et Sutherland~\cite{schwaber2020} dans le \textit{Scrum Guide}, ce framework repose sur trois piliers fondamentaux~: la \textbf{transparence}, l'\textbf{inspection} et l'\textbf{adaptation}.

\begin{figure}[H]
    \centering
    \fbox{\parbox{12cm}{\centering\vspace{2cm}\textit{Diagramme du processus Scrum}\vspace{2cm}}}
    \caption{Le processus Scrum~\cite{schwaber2020}}
    \label{fig:scrum_process}
\end{figure}

Le fonctionnement de Scrum s'organise autour de cycles itératifs appelés \textbf{sprints}, dont la durée est typiquement comprise entre une et quatre semaines~\cite{rubin2012}. Chaque sprint constitue une unité de travail autonome qui produit un \textbf{incrément fonctionnel} potentiellement livrable. Cette approche permet de recueillir un retour régulier des parties prenantes et d'ajuster les priorités en fonction de l'évolution des besoins.

\subsubsection{Les rôles Scrum}

Le framework Scrum définit trois rôles distincts dont les responsabilités sont clairement délimitées~\cite{schwaber2020}~:

\begin{itemize}[noitemsep]
    \item \textbf{Product Owner}~: responsable de la maximisation de la valeur du produit. Il gère et priorise le \textit{Product Backlog}, un artefact qui recense l'ensemble des fonctionnalités souhaitées~;
    \item \textbf{Scrum Master}~: garant du respect du cadre Scrum, il facilite les cérémonies et lève les obstacles rencontrés par l'équipe~;
    \item \textbf{Équipe de développement}~: un groupe auto-organisé et pluridisciplinaire chargé de la conception, du développement et des tests des incréments.
\end{itemize}

\subsubsection{Les événements Scrum}

Chaque sprint est ponctué par quatre événements formels~\cite{schwaber2020}~:

\begin{enumerate}[noitemsep]
    \item \textbf{Sprint Planning}~: réunion de planification où l'équipe sélectionne les éléments du backlog à réaliser durant le sprint~;
    \item \textbf{Daily Scrum}~: mêlée quotidienne de 15 minutes pour la synchronisation de l'équipe~;
    \item \textbf{Sprint Review}~: présentation de l'incrément fonctionnel aux parties prenantes à la fin du sprint~;
    \item \textbf{Sprint Retrospective}~: session de rétrospection pour identifier les améliorations du processus.
\end{enumerate}

\subsubsection{Les artefacts Scrum}

Le framework repose sur trois artefacts essentiels~\cite{rubin2012}~:

\begin{itemize}[noitemsep]
    \item \textbf{Product Backlog}~: liste ordonnée de l'ensemble des fonctionnalités, améliorations et corrections identifiées pour le produit~;
    \item \textbf{Sprint Backlog}~: sous-ensemble du Product Backlog sélectionné pour le sprint en cours, accompagné d'un plan de réalisation~;
    \item \textbf{Incrément}~: le résultat concret du sprint, constitué de l'ensemble des éléments du Product Backlog achevés, ajoutés aux incréments des sprints précédents.
\end{itemize}


\subsection{Justification du choix de Scrum}

Le choix de Scrum comme méthodologie de développement se justifie par plusieurs facteurs inhérents à la nature de notre projet~:

\begin{enumerate}[noitemsep]
    \item \textbf{Complexité du système}~: la plateforme LLM Dashboard comporte onze modules interconnectés. Scrum permet de décomposer cette complexité en sprints autonomes, chacun produisant un sous-ensemble fonctionnel cohérent~;
    \item \textbf{Besoins évolutifs}~: les exigences en matière de gouvernance des LLM évoluent rapidement, notamment avec l'apparition continue de nouveaux modèles et fournisseurs. L'approche itérative de Scrum permet d'intégrer ces évolutions au fil des sprints~;
    \item \textbf{Retour continu}~: les sprint reviews permettent de présenter régulièrement les fonctionnalités développées à l'entreprise d'accueil et d'ajuster les priorités~;
    \item \textbf{Taille de l'équipe}~: Scrum est particulièrement adapté aux petites équipes, ce qui correspond à la configuration de notre projet de fin d'études.
\end{enumerate}

\subsection{Application de Scrum au projet}

L'application de Scrum à notre projet s'est traduite par la définition de \textbf{trois sprints} couvrant l'intégralité du développement~:

\begin{table}[H]
\centering
\caption{Découpage du projet en sprints}
\label{tab:sprints}
\begin{tabularx}{\textwidth}{|c|l|X|}
\hline
\textbf{Sprint} & \textbf{Intitulé} & \textbf{Périmètre} \\
\hline
Sprint~1 & Infrastructure et Gateway LLM & Authentification, Gateway proxy, Gouvernance, Tableau de bord, base de données \\
\hline
Sprint~2 & Intelligence Artificielle & Détection d'anomalies, Explication IA, Prévision budgétaire, Analytics, Optimisation \\
\hline
Sprint~3 & Entreprise et Déploiement & Facturation, Traçage, Gestion des prompts, Tenants, Déploiement Docker, Tests \\
\hline
\end{tabularx}
\end{table}


% ────────────────────────────────────────────────────────────
\section{Méthodologies architecturales et d'ingénierie complémentaires}
\label{sec:methodologies_complementaires}
% ────────────────────────────────────────────────────────────

Si Scrum fournit le cadre organisationnel du développement itératif, la complexité technique de la plateforme LLM Dashboard nécessite l'adoption de méthodologies d'ingénierie complémentaires pour guider les choix architecturaux, opérationnels et liés à l'intelligence artificielle. Cette section présente les cinq approches retenues et leur rôle spécifique dans notre projet.

\subsection{Domain-Driven Design (DDD)}

Le \textbf{Domain-Driven Design}, introduit par Eric Evans~\cite{evans2003}, est une méthodologie architecturale centrée sur la modélisation fidèle du domaine métier dans le code source. Le principe fondamental consiste à aligner l'implémentation logicielle sur les concepts et la terminologie propres au domaine d'affaires, afin de produire un système maintenable et évolutif.

\textbf{Application au projet.} Notre plateforme couvre plusieurs domaines métier distincts~: l'\textit{économie des tokens} (pricing, agrégation, facturation), la \textit{gouvernance} (règles, quotas, restrictions), la \textit{détection d'anomalies} (baselines, scores ML) et l'\textit{identité multi-tenant} (utilisateurs, rôles, clés API). Le DDD nous a conduit à structurer le backend en \textbf{contextes délimités} (\textit{bounded contexts}), chacun encapsulé dans un module Python autonome avec ses propres modèles, services et repositories~:

\begin{itemize}[noitemsep]
    \item \texttt{modules/gateway/} --- Contexte \og Routage et proxy LLM \fg{}~;
    \item \texttt{modules/detection/} --- Contexte \og Analyse statistique et ML \fg{}~;
    \item \texttt{modules/billing/} --- Contexte \og Facturation contractuelle \fg{}.
\end{itemize}

Cette séparation empêche la complexité d'un domaine (par exemple la logique de calcul de factures) de polluer un autre domaine (la détection d'anomalies), et facilite l'ajout futur de nouveaux contextes à mesure que la réglementation de l'IA évolue.

\subsection{Platform Engineering}

Le \textbf{Platform Engineering} constitue l'évolution moderne du DevOps, centrée sur la création d'une \og plateforme interne de développement \fg{} (\textit{Internal Developer Platform --- IDP}) offrant un \og chemin doré \fg{} (\textit{golden path}) aux développeurs consommateurs~\cite{ford2017}.

\textbf{Application au projet.} Le LLM Dashboard agit précisément comme une \textbf{plateforme de services} pour les développeurs de QALITAS et GMAO PRO. Plutôt que de gérer eux-mêmes leurs clés API, leur monitoring et leur gouvernance, les équipes de développement de TIM Tunisie consomment un service centralisé et auto-administrable~:

\begin{itemize}[noitemsep]
    \item \textbf{Self-service}~: les administrateurs de tenants créent et gèrent leurs clés API, leurs règles de gouvernance et leurs contrats de facturation de manière autonome via le tableau de bord~;
    \item \textbf{Abstraction de la complexité}~: les agents applicatifs n'ont qu'un seul endpoint à consommer (\texttt{/v1/gateway/chat/completions}), indépendamment du fournisseur LLM sous-jacent~;
    \item \textbf{Observabilité intégrée}~: les métriques, les alertes et les rapports sont produits automatiquement, sans intervention des équipes consommatrices.
\end{itemize}


\subsection{Site Reliability Engineering (SRE)}

Le \textbf{Site Reliability Engineering}, formalisé par Google~\cite{beyer2016}, est une approche d'ingénierie logicielle appliquée aux opérations. Elle repose sur la définition de \textbf{Service Level Objectives} (SLO) et la gestion de budgets d'erreur, en opposition à l'objectif irréaliste de 100\,\% de disponibilité.

\textbf{Application au projet.} La plateforme LLM Dashboard joue un rôle critique dans la chaîne applicative de TIM Tunisie~: si le proxy Gateway est indisponible, \textbf{aucune application cliente ne peut accéder aux LLM}. Les principes SRE ont guidé plusieurs décisions architecturales~:

\begin{itemize}[noitemsep]
    \item \textbf{Métriques de latence}~: le Gateway mesure et journalise la latence de bout en bout (\textit{Time to First Token --- TTFT}) pour chaque appel~;
    \item \textbf{Health checks automatisés}~: chaque service Docker expose un endpoint \texttt{/health} interrogé par le scheduler~;
    \item \textbf{Circuit breaker}~: le mécanisme de fallback multi-fournisseur (Groq $\rightarrow$ OpenAI $\rightarrow$ Anthropic) garantit la résilience face aux pannes d'un fournisseur~;
    \item \textbf{Monitoring temps réel}~: les WebSockets du tableau de bord fournissent une visibilité instantanée sur les quatre signaux dorés (\textit{golden signals})~: latence, trafic, erreurs et saturation.
\end{itemize}


\subsection{MLOps}

Le \textbf{MLOps} (\textit{Machine Learning Operations}) intègre les pratiques DevOps à l'ingénierie du Machine Learning pour traiter les modèles comme des artefacts vivants nécessitant un monitoring et une mise à jour continus~\cite{sculley2015}.

\textbf{Application au projet.} La plateforme embarque trois algorithmes d'apprentissage automatique (STL, Isolation Forest, CUSUM) exécutés automatiquement par le pipeline Celery Beat. L'approche MLOps se manifeste dans~:

\begin{itemize}[noitemsep]
    \item \textbf{Monitoring continu des modèles}~: les baselines statistiques (\texttt{llm\_token\_baseline}) sont recalculées quotidiennement pour s'adapter à l'évolution naturelle des patterns de consommation, prévenant ainsi la dérive des données (\textit{data drift})~;
    \item \textbf{Pipeline automatisé}~: l'agrégation quotidienne, la détection d'anomalies et la génération de prévisions sont orchestrées par Celery Beat sans intervention humaine~;
    \item \textbf{Traçabilité}~: chaque anomalie détectée est persistée avec l'algorithme utilisé, le z-score calculé et les valeurs attendue/réelle, assurant la reproductibilité des résultats~;
    \item \textbf{Boucle de rétroaction}~: les explications IA générées par le module Explainer enrichissent la compréhension des anomalies et orientent l'ajustement des seuils de détection.
\end{itemize}


\subsection{GitOps}

Le \textbf{GitOps} est une méthodologie opérationnelle moderne dans laquelle le dépôt Git constitue la \og source de vérité unique \fg{} pour l'ensemble de l'infrastructure et de la configuration applicative~\cite{limoncelli2014}. Tout changement d'infrastructure ou de configuration est versionné, auditable et déployé de manière automatisée.

\textbf{Application au projet.} La gouvernance des LLM implique des configurations sensibles (règles de quota, limites budgétaires, restrictions de modèles) dont la traçabilité est essentielle. Notre approche GitOps se concrétise par~:

\begin{itemize}[noitemsep]
    \item \textbf{Infrastructure as Code}~: la configuration Docker Compose (6 services), les fichiers d'environnement et les scripts de migration sont versionnés dans le dépôt Git~;
    \item \textbf{CI/CD automatisé}~: le pipeline GitHub Actions exécute les tests, construit les images Docker et les déploie automatiquement~;
    \item \textbf{Piste d'audit}~: chaque modification de configuration est associée à un commit Git horodaté, ce qui est essentiel pour la conformité réglementaire du module de gouvernance.
\end{itemize}


\subsection{Synthèse comparative des méthodologies}

Le tableau~\ref{tab:methodologies_comparatif} résume les méthodologies adoptées et leur champ d'application dans notre projet.

\begin{table}[H]
\centering
\caption{Synthèse comparative des méthodologies adoptées}
\label{tab:methodologies_comparatif}
\begin{tabularx}{\textwidth}{|l|l|X|}
\hline
\textbf{Méthodologie} & \textbf{Focus principal} & \textbf{Application dans le projet} \\
\hline
Scrum & Vélocité et itération & Cadre organisationnel, livraison incrémentale UI/Backend \\
\hline
DDD & Complexité logique & Structure modulaire (Gateway, Billing, Detection) \\
\hline
SRE & Fiabilité et performance & Proxy Gateway critique, fallback multi-fournisseur \\
\hline
MLOps & Cycle de vie des modèles ML & Pipeline de détection d'anomalies, recalcul des baselines \\
\hline
Platform Eng. & Expérience développeur & Dashboard self-service pour les tenants \\
\hline
GitOps & Audit de configuration & Versionnement Docker, CI/CD, traçabilité gouvernance \\
\hline
\end{tabularx}
\end{table}


% ────────────────────────────────────────────────────────────
\section{Langage de modélisation~: UML}
\label{sec:uml}
% ────────────────────────────────────────────────────────────

Le \textbf{Unified Modeling Language} (UML) est le standard de modélisation graphique adopté pour la conception de notre système. Défini par l'Object Management Group (OMG), UML offre un ensemble de diagrammes complémentaires permettant de représenter la structure statique et le comportement dynamique d'un système logiciel~\cite{gamma1994}.

Dans le cadre de ce projet, nous utilisons les types de diagrammes suivants~:

\begin{itemize}[noitemsep]
    \item \textbf{Diagrammes de cas d'utilisation}~: pour identifier les acteurs et leurs interactions avec le système (global et par sprint)~;
    \item \textbf{Diagrammes de classes}~: pour modéliser la structure des entités métier et leurs relations~;
    \item \textbf{Diagrammes de séquence}~: pour détailler les flux d'interaction entre les composants du système~;
    \item \textbf{Diagrammes de composants}~: pour représenter l'architecture modulaire du backend~;
    \item \textbf{Diagrammes d'activité}~: pour illustrer les processus asynchrones (pipeline Celery)~;
    \item \textbf{Diagrammes d'états}~: pour modéliser les cycles de vie des entités (contrats, factures)~;
    \item \textbf{Diagrammes de déploiement}~: pour documenter l'architecture d'infrastructure Docker.
\end{itemize}

Les diagrammes ont été produits à l'aide de l'outil \textbf{PlantUML}, qui permet de générer des représentations UML à partir de descriptions textuelles versionnées dans le dépôt Git du projet.


\section*{Conclusion}

Ce premier chapitre a posé le cadre général de notre projet de fin d'études. La présentation de TIM Tunisie a mis en lumière le contexte industriel dans lequel s'inscrit notre travail, tandis que l'analyse de l'existant a révélé des lacunes critiques dans la gestion des LLM. La problématique formulée a guidé la définition d'une solution ambitieuse --- la plateforme LLM Dashboard --- dont les onze modules fonctionnels répondent de manière exhaustive aux besoins identifiés. La méthodologie Scrum, combinée aux approches architecturales complémentaires (DDD, SRE, MLOps, Platform Engineering, GitOps), garantit un développement à la fois agile, fiable et maintenable. Le chapitre suivant dresse un état de l'art des technologies fondamentales qui sous-tendent cette plateforme.
